\subsection{An actual reference algebra and the full simplicial prism}

\begin{theorem}[AT26: Proved; a coherent prism for the free-cell construction]
\label{AT26}
Work over $R=\mathbb Z_{(2)}$.
In AT25 assume in addition the actual global relative
free presentation
\[
 X=P\amalg_{\text{bottom }K(1)\text{ disks}}
          \operatorname{FreePD}(\text{the selected }D(2)\text{ disks}).
 \tag{AT.171}
\]
The bottom of the $i$th odd disk is the specified
$(dh_i,0)$; its top is $(x_i,y_i)$.
The algebra $P$ has $P_0=S$, coefficient $H^0$,
rational coefficient comparison, and in each simplex
is $S$ tensored over $R$ with free identity disks.
The maps $P\to A,B$ are specified.
This is a global free-cell presentation, not just
a description of its simplex-two algebra.

Let $Y$ be the simultaneous torsionfree resolved
extension in AT25 with its chosen target maps.
There are an actual simplicial strict PD algebra
$R_{\mathrm{ref}}$, maps
\[
 g:Y\longrightarrow R_{\mathrm{ref}},
 \qquad r_A:R_{\mathrm{ref}}\longrightarrow A,
 \qquad r_B:R_{\mathrm{ref}}\longrightarrow B,
 \tag{AT.172}
\]
and actual strict maps to the path objects
\[
 F_A:Y\longrightarrow A^{\Delta[1]},\qquad
 F_B:Y\longrightarrow B^{\Delta[1]},
 \qquad e_0F_C=r_Cg,\quad e_1F_C=f_C .
 \tag{AT.173}
\]
The reference has torsionfree terms, coefficient
$H^0$, rational coefficient comparison and integral
coefficient comparison through form three in every
simplex and weight.
Its support is $q\leq w$.
These assertions use the bounds of AT25.

The paths (AT.173) are specified on the whole
prism of every free parameter. No independently
chosen simplexwise maps are asserted to form a
homotopy.
\end{theorem}

\begin{proof}
Use the one common local standard model $G_N$
in AT20, where $N=2^e$ satisfies (AT.167).
For each selected cell use its actual elements
$h_i',x_i',y_i'$ in that model. Their identities are
\[
 \partial h_i'=\partial h_i,\qquad
 \partial x_i'=dh_i',\qquad dx_i'=y_i',
 \qquad \gamma_2(y_i')=x_i'y_i'=0.
 \tag{AT.174}
\]
These are the explicit divided-power formulas
of AT20, including $b_i\geq2N$.

Start with $P\otimes_S G_N$.
For each $i$ adjoin a relative even zero-form
$D(2)$ disk $(q_i,dq_i)$ whose bottom $K(1)$
value is
\[
                      (h_i-h_i',\,dh_i-dh_i').
 \tag{AT.175}
\]
The two endpoints of $h_i-h_i'$ are zero, so
these are actual bottom parameter maps.
Call the resulting algebra $R_{\mathrm{ref}}$.
In every simplex the new generator complements
are free even-to-odd identity disks.
They have their integral monomial contractions.
The $S$-linear coefficient contraction of $G_{N,n}$
and these disk contractions identify the vertical
cohomology of the reference with that of $P_n$.
The free identity-disk calculation in AT24
therefore gives the asserted cohomology through
degree three, coefficient $H^0$ and rational
comparison. All terms are free over $R$ and
all assigned weights satisfy the support bound.

The maps from $G_N$ to each target are the actual
coefficient-edge constructions of AT20.
Its element denoted $q_i$ in that proof is a
normalized two-simplex of the target with
boundary $f_C(h_i)-h_{i,C}'$.
Use precisely that element as the value of the
new reference generator in (AT.175).
The relative free universal property gives
$r_A,r_B$ on the whole reference algebra.

Define the map from $X$ by fixing its $P$ part
and setting
\[
             g(x_i)=x_i'+dq_i,\qquad
             g(y_i)=y_i'.
 \tag{AT.176}
\]
The zero face of the first expression is $dh_i$,
and its other faces vanish.
Its vertical differential is $y_i'$.
Thus (AT.176) respects the actual $D(2)$ bottom
data, rather than merely a formal top variable.
On the resolved $K(2)$ module belonging to $i$
set
\[
 g(u_0(t))=0,\qquad
 g(v_0(t))=-\varepsilon(t)q_i y_i',
 \qquad g(u_j)=g(v_j)=g(W^0)=0\quad(j>0).
 \tag{AT.177}
\]
The product $q_i y_i'$ has all faces zero, so
it is an actual $K(2)$ value.
Equations (AT.174) give
$d g(u_0)=\varepsilon\gamma_2(y_i')=0$ and
\[
 d g(v_0(t))=-\varepsilon(t)dq_i\,y_i'
       =-\varepsilon(t)(x_i'+dq_i)y_i'.
 \]
All other coefficient equations follow from
$\varepsilon\partial=0$, including the terminal
relation for $W^0$. This proves a strict map
on the actual resolved presentation.
It factors through the torsionfree image because
the reference is torsionfree.
This gives $g$ in (AT.172).

We now specify the complete prism.
Write $f_C$ for the chosen AT20 map on $Y$ and
$g_C=r_Cg$.
The carrier of AT24 gives the first top prism
value, denoted $H_0(e)$, on every normalized
two-dimensional parameter generator $e$.
It is an actual PD cochain map on the full
simplex-two algebra and has the face values
\[
        (d_0H_0(e),d_1H_0(e),d_2H_0(e),d_3H_0(e))
                         =(f_C(e),g_C(e),0,0).
 \tag{AT.178}
\]
For an old $D(2)$ generator its bottom is $a$,
where $f_C(a)=g_C(a)$.
For a resolved $K(2)$ generator put $a=0$.
Define the other two top prism values by
\[
                  H_1(e)=s_1g_C(e),\qquad
                  H_2(e)=s_2g_C(e).
 \tag{AT.179}
\]
Their four faces, in order, are
\[
 \begin{array}{c|cccc}
 H_1(e)&s_0f_C(a)&g_C(e)&g_C(e)&0\\
 H_2(e)&s_1f_C(a)&0&g_C(e)&g_C(e).
 \end{array}
 \tag{AT.180}
\]
On the bottom-edge prism take the two constant
values $s_0f_C(a)$ and $s_1f_C(a)$.
On all base-algebra simplices use the constant
homotopy of its specified map.

Here is an explicit completeness check for these
data. The pointed parameter for $D(2)$ is
$\Delta[2]/\Lambda^2_0$: its only non-base
nondegenerate simplices are the edge $a$ and
the triangle $e$. Its reduced free module has
normalization $R\xrightarrow{1}R$ in degrees
two and one. The parameter for $K(2)$ is
$\Delta[2]/\partial\Delta[2]$, with only $e$.
The triangular prism $\Delta[2]\times\Delta[1]$
has exactly three nondegenerate top tetrahedra,
indexed by the position where the interval
coordinate changes. They have the values
$H_0,H_1,H_2$ above.
With the standard ordered interval, its zero
endpoint is $g_C$ and its one endpoint is $f_C$:
the first top prism has $d_0=f_C$, and the last
has $d_3=g_C$.
Their common faces match by
(AT.178)--(AT.180); their exterior triangle
faces are the specified $f_C,g_C$, their
bottom-edge faces are the specified constant
edge prism, and their other boundary faces
are zero. The edge prism has precisely the
two indicated nondegenerate triangles.
Thus every nondegenerate simplex of each
parameter prism has a value, and all its
face equations hold. Degeneracies now define
the remaining values uniquely.

The free relative PD universal property turns
these parameter-prism maps into an algebra
homotopy. To check its vertical differential,
on the first tetrahedron use the actual
cochain map $H_0$; on the other two use the
cochain maps $s_1g_C,s_2g_C$.
The same check on the edge prism uses the
constant cochain map.
Every nonlinear twist and every $S$-linear
resolved relation is respected in each of
these actual algebra maps.
Their matching faces therefore respect the
relations on the entire prism.
Adjunction with the cotensor gives exactly
the strict maps (AT.173).

Equivalently, one may describe the universal
property degree by degree: the copower of
the source algebra by the finite set of
interval simplices is the coproduct of
those source copies. Its free parameters
are precisely the corresponding parameter
prisms. The assignments just checked give
compatible strict maps from each coproduct.
This proves an actual simplicial PD
homotopy, with no linear homotopy being
silently promoted to a multiplicative one.
\end{proof}
